% deepstudy-pdf template
% Copy this file to output/<slug>.tex and replace both @@PLACEHOLDER@@
% markers below. Don't use << >> for placeholders: Computer Modern
% ligates << and >> into « » guillemets when rendered, so they don't
% survive as literal text (verified — see SKILL.md Gotchas).
% Each %% TODO comment marks content the generating agent must write —
% real explanations, real code, real diagrams. Delete the TODO comments
% once filled in. Compile with compile.sh, never call pdflatex directly
% (minted needs -shell-escape and a 3-pass build for the TOC/index).
\documentclass[12pt,a4paper,oneside]{book}
\usepackage[utf8]{inputenc}
\usepackage[T1]{fontenc}
\usepackage{amsmath,amssymb}
\usepackage{listings}
\usepackage{xcolor}
\usepackage{hyperref}
\usepackage{tikz}
\usetikzlibrary{arrows.meta} % needed for modern arrow tips (Stealth, Latex, ...)
\usetikzlibrary{shapes.geometric} % needed for trapezium, diamond, etc.
\usetikzlibrary{positioning} % needed for "right=0.3cm of <node>" syntax
\usepackage{minted}
\usepackage{tcolorbox}
\usepackage{geometry}
\usepackage{fancyhdr}
\usepackage{imakeidx}

\makeindex[intoc]

\geometry{margin=2.5cm}
\setlength{\parskip}{0.8em}
\setlength{\parindent}{0pt}

\lstset{
    basicstyle=\ttfamily\small,
    breaklines=true,
    frame=single,
    numbers=left,
    numbersep=5pt,
    backgroundcolor=\color{gray!10},
    keywordstyle=\color{blue}\bfseries,
    commentstyle=\color{green!40!black},
    stringstyle=\color{red!40!black}
}

\pagestyle{fancy}
\fancyhf{}
\fancyhead[LE,RO]{\thepage}
\fancyhead[RE]{\nouppercase{\leftmark}}
\fancyhead[LO]{\nouppercase{\rightmark}}

\hypersetup{
    colorlinks=true,
    linkcolor=blue!60!black,
    urlcolor=blue!60!black,
    citecolor=blue!60!black,
}

% book class has no native \subtitle; fold it into \title as a second line
\title{@@TITULO@@\\[0.4em]\Large @@SUBTITULO@@}
\author{DeepStudy PDF Generator}
\date{\today}

\begin{document}

\maketitle
\tableofcontents

% CAPÍTULO 1: FUNDAMENTOS TEÓRICOS
\chapter{Fundamentos e Teoria Essencial}
%% TODO: explicar os conceitos base do assunto, por que ele existe,
%% que problema resolve, e um resumo da evolução histórica relevante.
%% Escreva de verdade — não deixe placeholder de lorem ipsum no PDF final.

% CAPÍTULO 2: ARQUITETURA INTERNA
\chapter{Por Debaixo dos Panos}
%% TODO: como o assunto funciona internamente. Use tikz para diagramas
%% de arquitetura reais (não decorativos) e, se houver código-fonte
%% relevante (ex.: de um projeto open source), analise trechos reais.

% CAPÍTULO 3: IMPLEMENTAÇÃO PRÁTICA
\chapter{Implementação e Exemplos}
%% TODO: múltiplos exemplos de código do básico ao avançado, usando
%% minted (```\begin{minted}{python}...\end{minted}```) para o
%% destaque de sintaxe real. Inclua benchmarks/análises de performance
%% quando fizer sentido para o assunto.

% CAPÍTULO 4: CASOS AVANÇADOS
\chapter{Tópicos Avançados}
%% TODO: edge cases, otimizações, debugging e troubleshooting reais —
%% coisas que só aparecem quando se usa o assunto em produção.

% APÊNDICES
\appendix
\chapter{Referências e Leitura Adicional}
%% TODO: fontes reais (docs oficiais, RFCs, papers, livros) — sem
%% inventar URLs.

\chapter{Glossário}
%% TODO: termos técnicos introduzidos nos capítulos anteriores.
%% Use \index{termo} nas ocorrências ao longo do texto para alimentar
%% o índice remissivo abaixo.

\printindex

\end{document}
